\chapter{Lower bounds of polynomials}
\label{chap:polybounds}

Every leaf test reduces to a statement of the form ``the polynomial $p$ is at least $m$
on $[-1,1]^5$'', where $p$ has integer coefficients and degree at most $2$ in each
variable (degree at most $4$ in $t_0, t_1$ for one family of orbit tests). This chapter
gives the two lower bounds that are used, and the integer encoding that lets the kernel
compute them quickly.

\begin{definition}
\label{def:pe}
\lean{Cayley.Pack.PE, Cayley.Pack.PE.ev, Cayley.Pack.PE.nrm}
\leanok
A \emph{polynomial expression} is a term built from integer constants, the variables
$t_0, \dots, t_4$, addition, subtraction and multiplication. It can be evaluated in any
commutative ring. Its \emph{norm} is computed by replacing every constant by its
absolute value and evaluating at $t = (1, \dots, 1)$; it bounds the sum of the absolute
values of the coefficients of the polynomial.
\end{definition}

The leaf tests are written once, generically in the arithmetic: the same definition,
instantiated with integers, is the checker; instantiated with polynomial expressions it
produces $p$; instantiated with reals it gives the geometric quantity. Maps that commute
with the ring operations commute with the generic definitions, which relates the three
instances.

\section{The coefficient bound}

\begin{theorem}[Coefficient bound]
\label{thm:coef_bound}
\lean{Cayley.Pack.lower_le_ev, Cayley.Pack.pLower_le, Cayley.Pack.pStats_sound}
\leanok
\uses{def:pe}
Let $p = \sum_m c_m t^m$. For $t \in [-1,1]^5$,
\[
  p(t) \ge c_0 - \sum_{m \text{ with an odd exponent}} |c_m| - \sum_{\substack{m \ne 0 \text{ all exponents even}\\ c_m < 0}} |c_m|,
  \qquad
  |p(t) - c_0| \le \sum_{m \ne 0} |c_m|.
\]
\end{theorem}
\begin{proof}
\leanok
For $|t_i| \le 1$ every monomial lies in $[-1, 1]$, and a monomial with only even
exponents lies in $[0, 1]$.
\end{proof}

\section{The Bernstein bound}

On $[-1, 1]$ a quadratic $a_0 + a_1 t + a_2 t^2$ equals
$b_0 \beta_0(t) + b_1 \beta_1(t) + b_2 \beta_2(t)$ with
\[
  (b_0, b_1, b_2) = (a_0 - a_1 + a_2,\ a_0 - a_2,\ a_0 + a_1 + a_2),\qquad
  \beta_0 = \Bigl(\frac{1 - t}{2}\Bigr)^2,\ \beta_1 = \frac{1 - t^2}{2},\ \beta_2 = \Bigl(\frac{1 + t}{2}\Bigr)^2.
\]
The $\beta_\kappa$ are nonnegative on $[-1, 1]$ and sum to $1$.

\begin{definition}
\label{def:bernstein}
\lean{Cayley.Pack.bT, Cayley.K.bern3}
\leanok
\uses{def:pe}
The \emph{Bernstein coefficients} of a polynomial of degree at most $2$ in each of five
variables are obtained by applying the transformation $(a_0, a_1, a_2) \mapsto (b_0, b_1, b_2)$
along each variable in turn. There are $3^5 = 243$ of them.
\end{definition}

\begin{theorem}[Bernstein bound]
\label{thm:bern_bound}
\lean{Cayley.Pack.bern_real, Cayley.Pack.bernGE_le}
\leanok
\uses{def:bernstein}
If every Bernstein coefficient of $p$ is at least $m$, then $p(t) \ge m$ for all
$t \in [-1,1]^5$.
\end{theorem}
\begin{proof}
\leanok
$p(t) = \sum_\kappa b_\kappa \prod_i \beta_{\kappa_i}(t_i)$, a convex combination of the
Bernstein coefficients.
\end{proof}

Neither bound dominates the other (for $p = t^2$ the coefficient bound is $0$ and the
Bernstein bound is $-1$), but on the gap polynomials the Bernstein bound is usually much
sharper; using it in addition halved the number of leaves in the certificate. It costs
more to compute, so the checker tries the coefficient bound first.

\section{Kronecker packing}

The kernel evaluates arithmetic on natural numbers with GMP, but each step of a
recursive function is comparatively expensive. The checker therefore represents a
polynomial by a single large integer.

\begin{definition}
\label{def:layout}
\lean{Cayley.Pack.Layout, Cayley.Pack.Layout.idx, Cayley.Pack.kron}
\leanok
\uses{def:pe}
A \emph{layout} assigns to every exponent vector $m$ with $m_i < r_i$ the slot
$\mathrm{idx}(m) = \sum_i m_i s_i$, where $s_i = \prod_{i' < i} r_{i'}$. For a slot width
$W$, the \emph{Kronecker point} is $t_i = 2^{W s_i}$.
\end{definition}

\begin{lemma}
\label{lem:kronecker}
\lean{Cayley.Pack.Layout.ev_kron, Cayley.Pack.enc_eq}
\leanok
\uses{def:layout}
If the exponents of $p$ fit the layout, then $p$ evaluated at the Kronecker point is
$\sum_j c_j 2^{W j}$, where $c_j$ is the coefficient in slot $j$. If moreover
$|c_j| < 2^{W-1}$ for all $j$, then adding $2^{W-1}$ to every slot gives a natural number
whose $j$-th base-$2^W$ digit is $c_j + 2^{W-1}$.
\end{lemma}
\begin{proof}
\leanok
Evaluation at the Kronecker point is a ring homomorphism, and the mixed-radix indexing
is injective on the exponents that fit.
\end{proof}

Since evaluation is a ring homomorphism, the checker computes $p$ at the Kronecker
point with a handful of big-integer additions and multiplications. Every coefficient
is then available as a digit: the top bit of each offset digit is its sign, a bitwise
\emph{and} with a mask selects a set of slots, and the sum of a set of digits is a
remainder modulo $2^W - 1$.

\begin{lemma}
\label{lem:packed_bounds}
\lean{Cayley.Pack.pLowerL_eq, Cayley.Pack.pStatsL_eq, Cayley.Pack.bern3_pv, Cayley.Pack.bernGE_spec}
\leanok
\uses{lem:kronecker, thm:coef_bound, def:bernstein}
If the coefficient norm of $p$ is below the bound fixed by the slot width, then the
checker's mask computations applied to $p$ at the Kronecker point return exactly the
coefficient bound of Theorem~\ref{thm:coef_bound}, and its packed Bernstein transform
returns the packed Bernstein coefficients. A Bernstein test that passes therefore
certifies that every Bernstein coefficient is at least $m$.
\end{lemma}
\begin{proof}
\leanok
\uses{lem:packN}
Digit-by-digit arithmetic on packed numbers (Lemma~\ref{lem:packN}). The Bernstein
transform along variable $i$ separates the slots by their digit $i$ with masks, shifts
them to a common position, and recombines them.
\end{proof}

\begin{lemma}
\label{lem:packN}
\lean{Cayley.Pack.packN, Cayley.Pack.packN_land, Cayley.Pack.packN_add, Cayley.Pack.packN_signs, Cayley.Pack.packN_mod}
\leanok
Let $\mathrm{pack}(d) = \sum_{j < n} d_j 2^{W j}$ with every digit $d_j < 2^W$. Bitwise
\emph{and}, digit-wise sums and differences (when no digit over- or underflows), the
extraction of the top bit of every digit, and the sum of all digits (via the remainder
modulo $2^W - 1$, when that sum is below $2^W - 1$) all act digit by digit.
\end{lemma}
\begin{proof}
\leanok
Induction on $n$.
\end{proof}

The checker uses slot width $W = 400$ and requires coefficient norms below $2^{390}$,
which leaves room for the Bernstein transform.
